\section{Концептуальна модель системи}

\subsection{Опис системи та процесу її функціонування}

Система, що моделюється, -- аеропорт, який приймає літаки на посадку. У ній можна виділити чотири підсистеми:
\begin{itemize}
\item \textbf{потік літаків} -- літаки, що потребують посадки, прибувають у випадкові моменти часу, у середньому 10~літаків за годину;
\item \textbf{посадкові смуги} -- $n$~однакових смуг; літак, що сів на смугу, звільняє її через 35~хвилин;
\item \textbf{черга очікування} -- літаки, які прибули, коли всі смуги зайняті, кружляють над аеропортом і чекають на посадку;
\item \textbf{фінансова підсистема} -- аеропорт отримує дохід за кожну посадку, сплачує за дозаправку літаків, що довго чекають, і вкладає кошти в будівництво додаткових смуг.
\end{itemize}

Процес обслуговування одного літака показано на рисунку~\ref{fig:process}. Якщо в момент прибуття є вільна смуга, літак одразу сідає, і аеропорт отримує 2000~грн. Інакше літак стає в чергу. З моменту прибуття для нього відраховуються два інтервали. Перший -- поріг дозаправки, від 60 до 80~хвилин: якщо літак досі не сів, аеропорт дозаправляє його, витрачаючи від 800 до 1200~грн. Другий -- 140~хвилин: якщо літак за цей час не сів, він відлітає в інший аеропорт, і аеропорт нічого за нього не отримує. Літак, що дочекався вільної смуги, сідає й приносить від 1400 до 1600~грн.

\begin{figure}[H]
\centering
\begin{tikzpicture}[
  font=\fontsize{11}{13}\selectfont,
  node distance=0.7cm and 1.0cm,
  box/.style={draw, rounded corners=2pt, align=center, minimum height=1.0cm, text width=3.9cm},
  money/.style={box, fill=gray!12},
  cond/.style={draw, diamond, aspect=2.4, align=center, inner sep=1pt, text width=2.9cm},
  arr/.style={-{Stealth[length=2.2mm]}},
]
\node[box] (arrive) {Прибуття літака\\(у середньому 10 за год)};
\node[cond, below=of arrive] (free) {Є вільна\\смуга?};
\node[money, right=1.6cm of free] (now) {Посадка без очікування\\дохід 2000 грн};
\node[box, below=of free] (queue) {Очікування в черзі};
\node[cond, below=of queue] (got) {Смуга звільнилась\\до 140 хв?};
\node[money, right=1.6cm of got] (late) {Посадка з очікуванням\\дохід $1500 \pm 100$ грн};
\node[box, below=of got] (leave) {Відліт в інший аеропорт\\дохід 0 грн};
\node[money, left=1.2cm of queue] (fuel) {Дозаправка, якщо чекає\\довше $70 \pm 10$ хв\\витрати $1000 \pm 200$ грн};
\node[box] (busy) at (now |- queue) {Смуга зайнята 35 хв,\\потім вільна};
\draw[arr] (arrive) -- (free);
\draw[arr] (free) -- node[above]{так} (now);
\draw[arr] (free) -- node[right]{ні} (queue);
\draw[arr] (queue) -- (got);
\draw[arr] (got) -- node[above]{так} (late);
\draw[arr] (got) -- node[right]{ні} (leave);
\draw[arr, dashed] (queue) -- (fuel);
\draw[arr] (now) -- (busy);
\draw[arr] (late) -- (busy);
\end{tikzpicture}
\caption{Процес обслуговування літака в аеропорту}
\label{fig:process}
\end{figure}

Літаки в черзі сідають у порядку прибуття. Коли смуга звільняється, першим сідає літак, що чекає найдовше, а не той, що прибув саме в цей момент. Дозаправка не зупиняє відлік до відльоту: дозаправлений літак, що так і не дочекався смуги, через 140~хвилин після прибуття все одно відлітає, і витрати на його дозаправку аеропорт уже поніс.

Аеропорт може збудувати від 0 до 10~додаткових смуг, кожна коштує 3\,000\,000~грн. Додаткові смуги прискорюють посадку, скорочують чергу, кількість дозаправок і відльотів, а отже підвищують річний прибуток. Час окупності $T = S/p$ показує, за скільки часу цей приріст прибутку поверне витрати $S$ на будівництво. Метою моделювання є кількість смуг, за якої $T$ мінімальний.

\subsection{Припущення моделі}

Текст завдання залишає кілька моментів невизначеними. У моделі прийнято такі припущення:
\begin{itemize}
\item запис «$70 \pm 10$ хвилин» означає, що поріг дозаправки кожного літака -- рівномірно розподілена величина на $[60; 80]$~хв; так само дохід $1500 \pm 100$~грн і витрати $1000 \pm 200$~грн рівномірно розподілені на $[1400; 1600]$ і $[800; 1200]$~грн;
\item літак дозаправляється не більше одного разу, а поріг дозаправки (не більше 80~хв) завжди менший за час відльоту (140~хв);
\item смуга зайнята 35~хвилин незалежно від того, чекав літак чи ні; смуги не відмовляють і однакові;
\item черга не обмежена за довжиною, літаки в ній обслуговуються в порядку прибуття;
\item у прибуток $p$ входять лише дохід від посадок і витрати на дозаправку, перераховані на 365~діб; інших витрат аеропорту в завданні немає;
\item $S$ -- витрати лише на додаткові смуги, тобто $S = 3 \cdot 10^6 \cdot k$~грн за $k$ додаткових смуг.
\end{itemize}

Прибуток $p$ у формулі окупності можна розуміти двояко. Окупність інвестицій зазвичай рахують за тим прибутком, який ці інвестиції приносять, тобто за приростом річного прибутку порівняно з аеропортом без нових смуг. У роботі це основна інтерпретація. Друга -- загальний річний прибуток аеропорту з новими смугами; за двох смуг він від'ємний, тож вона має сенс лише за достатньої кількості смуг. Модель обчислює обидва варіанти, а розділ~\ref{sec:experiment} показує, що оптимум за ними однаковий.

\subsection{Вхідні змінні моделі}

Вхідні змінні та параметри моделі наведено в таблиці~\ref{tab:inputs}. Керованою змінною є лише кількість посадкових смуг~$n$; решта параметрів задані умовою завдання.

\begin{kpitable}{|L{5.6cm}|L{2.7cm}|L{4.1cm}|L{3.0cm}|}{Вхідні змінні моделі\label{tab:inputs}}{\textbf{Змінна} & \textbf{Позначення} & \textbf{Значення} & \textbf{Одиниці}}
Кількість посадкових смуг (керована) & $n$ & $2 + k$, $k = 0, \ldots, 10$ & смуг \\ \hline
Інтенсивність надходження літаків & $\lambda$ & 10, експоненційні інтервали із середнім 6~хв & літаків / год \\ \hline
Час зайнятості смуги після посадки & $t_s$ & 35 & хв \\ \hline
Поріг очікування до дозаправки & $t_r$ & рівномірно на $[60; 80]$ & хв \\ \hline
Час очікування до відльоту & $t_l$ & 140 & хв \\ \hline
Дохід за посадку без очікування & $c_0$ & 2000 & грн \\ \hline
Дохід за посадку з очікуванням & $c_1$ & рівномірно на $[1400; 1600]$ & грн \\ \hline
Витрати на дозаправку & $z$ & рівномірно на $[800; 1200]$ & грн \\ \hline
Вартість однієї додаткової смуги & $s$ & 3\,000\,000 & грн \\ \hline
Найбільша кількість додаткових смуг & $k_{\max}$ & 10 & смуг \\ \hline
\end{kpitable}

\subsection{Вихідні змінні моделі}

Вихідні змінні наведено в таблиці~\ref{tab:outputs}. Цільовою характеристикою є час окупності $T$; решта пояснюють, звідки береться прибуток, і потрібні для верифікації моделі.

\begin{kpitable}{|L{7.0cm}|L{2.7cm}|L{5.7cm}|}{Вихідні змінні моделі\label{tab:outputs}}{\textbf{Змінна} & \textbf{Позначення} & \textbf{Одиниці}}
Кількість літаків, що прибули & $N_a$ & літаків \\ \hline
Кількість посадок без очікування & $N_0$ & літаків \\ \hline
Кількість посадок з очікуванням & $N_1$ & літаків \\ \hline
Кількість дозаправок & $N_r$ & літаків \\ \hline
Кількість літаків, що відлетіли в інший аеропорт & $N_l$ & літаків \\ \hline
Середній час очікування посадки & $W$ & хв \\ \hline
Середня та найбільша довжина черги & $L$, $L_{\max}$ & літаків \\ \hline
Середня кількість зайнятих смуг & $B$ & смуг \\ \hline
Дохід від посадок і витрати на дозаправку & $D$, $C$ & грн \\ \hline
Річний прибуток аеропорту & $P$ & грн / рік \\ \hline
Приріст річного прибутку порівняно з $n = 2$ & $\Delta P$ & грн / рік \\ \hline
Час окупності додаткових смуг & $T$ & місяців \\ \hline
\end{kpitable}

\subsection{Висновки до розділу}

Аеропорт описано як систему з потоком літаків, $n$~посадковими смугами, чергою очікування і фінансовою підсистемою. Особливість системи -- нетерплячі літаки з двома індивідуальними порогами очікування: після $70 \pm 10$~хв літак дозаправляють, після 140~хв він відлітає. Визначено вхідні змінні (керована -- кількість смуг $n$ від 2 до 12) і вихідні змінні, серед яких цільова -- час окупності додаткових смуг $T$. Сформульовано припущення, що уточнюють текст завдання.
